\documentclass[11pt,a4paper]{article}

% Required packages
\usepackage{graphicx}
\usepackage{float}
\usepackage{caption}
\usepackage{adjustbox}

\captionsetup{
    font=small,
    labelfont=bf
}

\usepackage[margin=1.5cm]{geometry}
\usepackage{pdflscape}
\usepackage{amsmath}
\usepackage{amsfonts}

\title{Complete Table of Contents\\Chapters 1--4\\[0.3cm]\small{Randomized Reward Shaping Deep Reinforcement Learning\\for Self-Adaptive Systems}}
\author{Alireza Kavianifar}
\date{\today}

\begin{document}

\maketitle

\tableofcontents
\newpage

% ============================================================================
% CHAPTER 1: INTRODUCTION
% ============================================================================

\setcounter{section}{0}
\renewcommand{\thesection}{1.\arabic{section}}
\renewcommand{\thesubsection}{1.\arabic{section}.\arabic{subsection}}
\renewcommand{\thesubsubsection}{1.\arabic{section}.\arabic{subsection}.\arabic{subsubsection}}

\section{Background and Context}

This section establishes the foundational context for self-adaptive systems and their role in managing runtime uncertainty.

\subsection{Self-Adaptive Systems: motivation and examples}

Self-adaptive systems autonomously adjust their behavior in response to changing environmental conditions and internal states. Examples include IoT networks, cloud computing infrastructures, and autonomous vehicles that must maintain quality-of-service goals despite unpredictable operating conditions.

\subsection{The MAPE-K feedback loop and the Managed/Managing System abstraction}

The MAPE-K (Monitor, Analyze, Plan, Execute, Knowledge) feedback loop provides a reference architecture for self-adaptive systems. The Managed System executes application logic, while the Managing System implements adaptation logic through continuous monitoring, analysis, planning, and execution cycles.

\subsection{Limitations of rule-based and static adaptation}

Traditional rule-based adaptation approaches struggle with:
\begin{itemize}
    \item Large adaptation spaces where manual rule definition becomes intractable
    \item Dynamic environments where pre-defined rules cannot anticipate all scenarios
    \item Complex trade-offs between multiple quality-of-service objectives
    \item Inability to learn from experience and improve over time
\end{itemize}

\section{Runtime Uncertainty in Self-Adaptive Systems}

Uncertainty represents a fundamental challenge for adaptation decision-making, requiring systems to make effective choices despite incomplete or imperfect information.

\subsection{Sources of uncertainty (environmental, system-level)}

Uncertainty arises from multiple sources:
\begin{itemize}
    \item \textbf{Environmental uncertainty}: Unpredictable interference, network conditions, workload variations
    \item \textbf{System-level uncertainty}: Hardware failures, software bugs, resource contention
    \item \textbf{Model uncertainty}: Imperfect models of system behavior and environment dynamics
    \item \textbf{Goal uncertainty}: Evolving or conflicting adaptation objectives
\end{itemize}

\subsection{Impact of uncertainty on adaptation decisions (examples from DeltaIoT)}

Using the DeltaIoT case study, we demonstrate how uncertainty affects adaptation outcomes. For example, network interference variations can cause the same configuration to produce different quality-of-service results, making it difficult to predict which adaptation option will achieve goals.

\section{Learning-based Adaptation: Why Reinforcement Learning?}

Reinforcement learning offers a principled framework for learning adaptation policies through trial-and-error interaction with uncertain environments.

\subsection{RL benefits for runtime decision-making}

Reinforcement learning provides:
\begin{itemize}
    \item Automatic policy learning without manual rule engineering
    \item Ability to balance exploration and exploitation
    \item Handling of delayed rewards and long-term consequences
    \item Adaptation to changing environment dynamics
\end{itemize}

\subsection{Limits of classic RL in large adaptation spaces}

Classic RL approaches face challenges in self-adaptive systems:
\begin{itemize}
    \item Sample inefficiency in large discrete action spaces
    \item Difficulty generalizing across heterogeneous environments
    \item Convergence to local optima
    \item High computational cost of retraining
\end{itemize}

\section{What ``Continuous'' Means in This Thesis}

Clarifying terminology is essential to avoid confusion with continuous-action reinforcement learning.

\subsection{Continuous = continuous runtime learning / streaming adaptation cycles (not continuous-action control)}

In this thesis, ``continuous'' refers to:
\begin{itemize}
    \item \textbf{Continuous runtime learning}: Ongoing policy improvement through streaming adaptation cycles
    \item \textbf{Continuous operation}: Uninterrupted system operation during learning
    \item \textbf{NOT continuous-action spaces}: We use discrete adaptation options (DQN), not continuous control (PPO/DDPG)
\end{itemize}

\section{Problem Statement (precise formulation extracted from the article)}

\textbf{Research Problem}: How can self-adaptive systems effectively learn adaptation policies in large discrete adaptation spaces under runtime uncertainty, while maintaining continuous operation and achieving rapid convergence across heterogeneous environments?

The problem encompasses:
\begin{itemize}
    \item Large adaptation spaces (hundreds to thousands of discrete options)
    \item Runtime uncertainty from environmental and system-level sources
    \item Need for sample-efficient learning
    \item Requirement for policy generalization across diverse conditions
    \item Constraint of continuous system operation during learning
\end{itemize}

\section{Proposed Solution Overview: RS-DRL (high-level)}

We propose Randomized Reward Shaping Deep Reinforcement Learning (RS-DRL), a novel approach that enhances DQN with selective reward reshaping to accelerate learning and improve generalization.

\subsection{Key idea: randomized reward reshaping during experience replay}

The core innovation is a randomized reward reshaping mechanism that:
\begin{itemize}
    \item Selectively reshapes failed experiences during replay
    \item Treats certain failures as optimistic successes with probability $\rho$
    \item Accelerates learning by providing additional positive signals
    \item Promotes policy generalization across heterogeneous environments
    \item Operates during training without affecting runtime decision-making
\end{itemize}

\subsection{Architectural placement within MAPE-K (on-demand retraining \& async learning)}

RS-DRL integrates with MAPE-K through:
\begin{itemize}
    \item \textbf{Planning phase enhancement}: Learned policies guide adaptation option selection
    \item \textbf{Asynchronous learning}: Policy training occurs independently of runtime adaptation
    \item \textbf{On-demand retraining}: Triggered when environment changes are detected
    \item \textbf{Knowledge repository}: Stores trained models for runtime use
\end{itemize}

\section{Research Objectives and Research Questions}

\textbf{Primary Objective}: Develop and validate a learning-based approach for self-adaptive systems that achieves effective adaptation in large discrete spaces under uncertainty.

\textbf{Research Questions}:
\begin{enumerate}
    \item How can reward shaping accelerate learning in large discrete adaptation spaces?
    \item What mechanisms enable policy generalization across heterogeneous environments?
    \item How does RS-DRL compare to baseline DQN and state-of-the-art approaches?
    \item What architectural patterns support integration of learning with MAPE-K?
\end{enumerate}

\section{Contributions of the Thesis (expanded from the article's contribution list)}

This thesis makes the following contributions:

\begin{enumerate}
    \item \textbf{Novel Algorithm}: RS-DRL with randomized reward reshaping mechanism
    \item \textbf{Architectural Integration}: Process model for integrating RS-DRL with MAPE-K
    \item \textbf{Empirical Validation}: Comprehensive evaluation on DeltaIoT case study
    \item \textbf{Generalization Analysis}: Demonstration of cross-environment policy transfer
    \item \textbf{Comparative Study}: Benchmarking against DQN, HER, and search-based methods
\end{enumerate}

\section{Research Methodology Overview}

The research follows a systematic methodology combining theoretical modeling, algorithm design, and empirical validation.

\subsection{Theoretical modeling and RL formulation}

\begin{itemize}
    \item Formalize self-adaptive system as Markov Decision Process
    \item Define state space, action space, and reward functions
    \item Establish theoretical foundations for reward shaping
    \item Analyze convergence properties
\end{itemize}

\subsection{Algorithm design and reward shaping experiments}

\begin{itemize}
    \item Design RS-DRL algorithm with randomized reward reshaping
    \item Implement using Stable-Baselines3 DQN framework
    \item Conduct ablation studies on reshaping probability $\rho$
    \item Optimize hyperparameters through systematic experimentation
\end{itemize}

\subsection{Case study (DeltaIoT): datasets, training regime, evaluation metrics}

\begin{itemize}
    \item Use DeltaIoT IoT network as motivating case study
    \item Train on datasets from Gheibi \& Weyns with varying uncertainty
    \item Evaluate using metrics: cumulative reward, convergence speed, generalization
    \item Compare against baseline and state-of-the-art approaches
\end{itemize}

\section{Thesis Organization (brief outline of subsequent chapters)}

The remainder of this thesis is organized as follows:

\begin{itemize}
    \item \textbf{Chapter 2}: Foundations and related work on self-adaptive systems, RL, and reward shaping
    \item \textbf{Chapter 3}: System model and problem formulation using DeltaIoT case study
    \item \textbf{Chapter 4}: Architecture and RS-DRL method with process model
    \item \textbf{Chapter 5}: Experimental evaluation and results analysis
    \item \textbf{Chapter 6}: Discussion, limitations, and future work
    \item \textbf{Chapter 7}: Conclusions and contributions summary
\end{itemize}

\newpage

% ============================================================================
% CHAPTER 2: FOUNDATIONS AND RELATED WORK
% ============================================================================

\setcounter{section}{0}
\renewcommand{\thesection}{2.\arabic{section}}
\renewcommand{\thesubsection}{2.\arabic{section}.\arabic{subsection}}
\renewcommand{\thesubsubsection}{2.\arabic{section}.\arabic{subsection}.\arabic{subsubsection}}

\section{Self-Adaptive Systems: Concepts and Architectural Patterns}

This section establishes the foundational concepts of self-adaptive systems and the MAPE-K reference architecture.

\subsection{MAPE-K detailed (Monitor, Analyze, Plan, Execute, Knowledge)}

The MAPE-K feedback loop consists of five components:

\begin{itemize}
    \item \textbf{Monitor}: Collects runtime data from managed system and environment
    \item \textbf{Analyze}: Evaluates current state against adaptation goals
    \item \textbf{Plan}: Selects adaptation actions to achieve goals
    \item \textbf{Execute}: Applies selected adaptations to managed system
    \item \textbf{Knowledge}: Stores models, policies, and system state history
\end{itemize}

Each component has well-defined responsibilities and interfaces, enabling modular design and systematic reasoning about adaptation behavior.

\subsection{Managed vs Managing System abstraction and temporal classification (monitoring vs asynchronous learning)}

The Managed/Managing System abstraction separates:

\begin{itemize}
    \item \textbf{Managed System}: Executes application logic, provides observability
    \item \textbf{Managing System}: Implements adaptation logic through MAPE-K
\end{itemize}

Temporal classification distinguishes:
\begin{itemize}
    \item \textbf{Continuous processes}: Monitoring (always active)
    \item \textbf{Event-driven processes}: Analysis, Planning, Execution (triggered by conditions)
    \item \textbf{Asynchronous processes}: Learning (decoupled from runtime adaptation)
\end{itemize}

\section{Uncertainty in Self-Adaptive Systems}

Uncertainty is pervasive in self-adaptive systems and fundamentally affects adaptation decision-making.

\subsection{Taxonomy of uncertainties (environmental, systemic, goal drift)}

We adopt a taxonomy that categorizes uncertainty sources:

\begin{itemize}
    \item \textbf{Environmental uncertainty}: External factors (interference, workload, user behavior)
    \item \textbf{Systemic uncertainty}: Internal factors (hardware variability, software bugs)
    \item \textbf{Model uncertainty}: Imperfect models of system and environment
    \item \textbf{Goal uncertainty}: Evolving or conflicting objectives
\end{itemize}

\subsection{Effects of uncertainty on verification and adaptation (motivating evidence from DeltaIoT)}

Using DeltaIoT as evidence, we demonstrate:

\begin{itemize}
    \item Same configuration produces different QoS under varying interference
    \item Verification results show probability distributions, not deterministic outcomes
    \item Adaptation decisions must account for uncertainty to avoid goal violations
    \item Learning-based approaches can adapt to uncertainty through experience
\end{itemize}

Figure 2 from the article shows how packet loss varies across interference levels, motivating the need for uncertainty-aware adaptation.

\section{Reinforcement Learning Foundations for SAS}

This section establishes RL foundations relevant to self-adaptive systems.

\subsection{MDPs, policy/value formulation and why RL suits adaptation problems}

\textbf{Markov Decision Process (MDP)} formulation:
\begin{itemize}
    \item State space $\mathcal{S}$: System and environment observations
    \item Action space $\mathcal{A}$: Available adaptation options
    \item Transition function $P(s'|s,a)$: State dynamics
    \item Reward function $R(s,a,s')$: Quality-of-service feedback
    \item Discount factor $\gamma$: Future reward weighting
\end{itemize}

\textbf{Policy and value functions}:
\begin{itemize}
    \item Policy $\pi(a|s)$: Probability of selecting action $a$ in state $s$
    \item State-value $V^\pi(s)$: Expected return from state $s$ under policy $\pi$
    \item Action-value $Q^\pi(s,a)$: Expected return from taking action $a$ in state $s$
\end{itemize}

RL suits adaptation because it:
\begin{itemize}
    \item Handles sequential decision-making under uncertainty
    \item Learns from experience without requiring complete system models
    \item Balances immediate and long-term objectives through discount factor
    \item Adapts to changing environment dynamics
\end{itemize}

\subsection{Known challenges: local optima, generalization, retraining cost}

Classic RL faces challenges in self-adaptive systems:

\begin{itemize}
    \item \textbf{Local optima}: Greedy policies may converge to suboptimal solutions
    \item \textbf{Sample inefficiency}: Large action spaces require many experiences
    \item \textbf{Generalization}: Policies may not transfer across environments
    \item \textbf{Retraining cost}: Computational overhead of policy updates
    \item \textbf{Exploration-exploitation}: Balancing learning and performance
\end{itemize}

\section{Deep Reinforcement Learning in Large Discrete Adaptation Spaces}

Deep RL extends classic RL to handle large state and action spaces through function approximation.

\subsection{DQN architecture and rationale for discrete adaptation options (why DQN was chosen in the article)}

\textbf{Deep Q-Network (DQN)} architecture:
\begin{itemize}
    \item Neural network approximates $Q(s,a)$ function
    \item Input: State representation $s \in \mathbb{R}^d$
    \item Output: Q-values for all actions $[Q(s,a_1), \ldots, Q(s,a_n)]$
    \item Loss: Temporal difference error minimization
\end{itemize}

\textbf{Why DQN for self-adaptive systems}:
\begin{itemize}
    \item Adaptation spaces are naturally discrete (configuration options)
    \item DQN handles large discrete action spaces efficiently
    \item Single forward pass evaluates all options simultaneously
    \item Proven stability and convergence properties
\end{itemize}

\subsection{Experience replay, target networks, stability techniques used}

DQN employs several techniques for stable learning:

\begin{itemize}
    \item \textbf{Experience replay}: Store experiences $(s,a,r,s')$ in memory, sample randomly for training
    \item \textbf{Target network}: Separate network for computing TD targets, updated periodically
    \item \textbf{$\epsilon$-greedy exploration}: Balance exploration and exploitation
    \item \textbf{Gradient clipping}: Prevent exploding gradients
    \item \textbf{Reward normalization}: Stabilize learning across different scales
\end{itemize}

\section{Reward Shaping and Experience Relabeling (state of the art)}

Reward shaping techniques modify reward signals to accelerate learning and improve generalization.

\subsection{HER and goal-conditioned methods --- contrast to RS-DRL (no pre-defined goals)}

\textbf{Hindsight Experience Replay (HER)}:
\begin{itemize}
    \item Relabels failed experiences with achieved goals
    \item Requires explicit goal representation in state space
    \item Effective for sparse reward environments
    \item Assumes goals can be extracted from final states
\end{itemize}

\textbf{Contrast with RS-DRL}:
\begin{itemize}
    \item RS-DRL does not require pre-defined goal representations
    \item Reshaping based on reward thresholds, not goal achievement
    \item Applicable to multi-objective optimization without explicit goals
    \item Randomized reshaping provides stochastic exploration benefit
\end{itemize}

\subsection{Other generalization techniques}

Additional techniques for improving RL generalization:
\begin{itemize}
    \item \textbf{Domain randomization}: Train on diverse environment variations
    \item \textbf{Meta-learning}: Learn to adapt quickly to new environments
    \item \textbf{Transfer learning}: Initialize policies from related tasks
    \item \textbf{Curriculum learning}: Gradually increase task difficulty
\end{itemize}

\section{Related Work: classification and gaps}

This section surveys existing approaches to learning-based adaptation and identifies research gaps.

\subsection{ML-based methods, search-based planners, RL-based adaptation --- strengths and weaknesses (article's survey)}

\textbf{Machine Learning-based methods}:
\begin{itemize}
    \item Supervised learning for adaptation prediction
    \item Strengths: Fast inference, interpretable models
    \item Weaknesses: Require labeled training data, limited to seen scenarios
\end{itemize}

\textbf{Search-based planners}:
\begin{itemize}
    \item Optimization algorithms (genetic algorithms, simulated annealing)
    \item Strengths: No training required, handle complex objectives
    \item Weaknesses: Computational overhead, may not scale to large spaces
\end{itemize}

\textbf{RL-based adaptation}:
\begin{itemize}
    \item Learn policies through trial-and-error
    \item Strengths: Adapt to uncertainty, improve over time
    \item Weaknesses: Sample inefficiency, generalization challenges
\end{itemize}

\subsection{Identified gap: scalable, uncertainty-aware, continuous runtime DRL for large adaptation spaces}

\textbf{Research gap}: Existing approaches lack:
\begin{itemize}
    \item Scalable learning in large discrete adaptation spaces (100s-1000s options)
    \item Effective generalization across heterogeneous environments
    \item Sample-efficient learning under runtime uncertainty
    \item Continuous operation during learning (asynchronous training)
    \item Principled reward shaping for multi-objective adaptation
\end{itemize}

RS-DRL addresses this gap through randomized reward reshaping and asynchronous learning architecture.

\section{Chapter Summary (lead into Chapter 3's formal modeling)}

This chapter established foundational concepts:
\begin{itemize}
    \item Self-adaptive systems and MAPE-K architecture
    \item Uncertainty taxonomy and impact on adaptation
    \item Reinforcement learning formulation and DQN architecture
    \item Reward shaping techniques and related work
    \item Research gap motivating RS-DRL approach
\end{itemize}

Chapter 3 formalizes the adaptation problem using the DeltaIoT case study, providing concrete MDP formulation and experimental setup.

\newpage

% ============================================================================
% CHAPTER 3: SYSTEM MODEL AND PROBLEM FORMULATION
% ============================================================================

\setcounter{section}{0}
\renewcommand{\thesection}{3.\arabic{section}}
\renewcommand{\thesubsection}{3.\arabic{section}.\arabic{subsection}}
\renewcommand{\thesubsubsection}{3.\arabic{section}.\arabic{subsection}.\arabic{subsubsection}}

\section{Motivating Case Study: DeltaIoT System (detailed description)}

DeltaIoT serves as the motivating case study, providing a concrete context for evaluating RS-DRL.

\subsection{Network, motes, links, configuration parameters and examples}

\textbf{DeltaIoT system architecture}:
\begin{itemize}
    \item IoT sensor network with multiple motes (sensor nodes)
    \item Wireless communication links subject to interference
    \item Gateway collecting sensor data
    \item Quality-of-service requirements: packet loss, latency, energy consumption
\end{itemize}

\textbf{Configuration parameters}:
\begin{itemize}
    \item Transmission power levels (affects range and energy)
    \item Distribution strategy (single-path vs multi-path routing)
    \item Packet size (affects transmission time and reliability)
\end{itemize}

\textbf{Example configurations}:
\begin{itemize}
    \item Low power + single-path: Energy-efficient but vulnerable to interference
    \item High power + multi-path: Reliable but energy-intensive
    \item Medium power + adaptive routing: Balanced trade-off
\end{itemize}

\section{Modeling Uncertainty in DeltaIoT and SAS}

Uncertainty in DeltaIoT arises from environmental interference and message load variability.

\subsection{Interference/noise ranges, message-load variability, data shifts (Fig.2 evidence)}

\textbf{Interference uncertainty}:
\begin{itemize}
    \item Interference levels vary from 0.0 to 1.0 (normalized)
    \item Affects packet loss probability and transmission reliability
    \item Figure 2 shows packet loss increasing with interference
    \item Same configuration produces different outcomes under different interference
\end{itemize}

\textbf{Message load variability}:
\begin{itemize}
    \item Number of messages varies over time
    \item Affects network congestion and latency
    \item Creates temporal dependencies in adaptation decisions
\end{itemize}

\textbf{Data shifts}:
\begin{itemize}
    \item Training and deployment environments may differ
    \item Requires policy generalization across interference ranges
    \item Motivates need for robust learning approaches
\end{itemize}

\subsection{How uncertainty affects QoS metrics (packet loss, latency, energy)}

\textbf{Packet loss}:
\begin{itemize}
    \item Increases with interference and low transmission power
    \item Probabilistic: same configuration yields different loss rates
    \item Critical metric for reliability requirements
\end{itemize}

\textbf{Latency}:
\begin{itemize}
    \item Affected by retransmissions due to packet loss
    \item Increases with network congestion
    \item Important for time-sensitive applications
\end{itemize}

\textbf{Energy consumption}:
\begin{itemize}
    \item Increases with transmission power and retransmissions
    \item Trade-off with reliability and latency
    \item Critical for battery-powered sensor nodes
\end{itemize}

\section{Adaptation Space and Configuration Representation}

The adaptation space defines the set of available configuration options.

\subsection{Definition: adaptation option, parameter encoding, scale (216 / 1,296 / 4,096)}

\textbf{Adaptation option}: A complete system configuration specifying all parameter values.

\textbf{Parameter encoding}:
\begin{itemize}
    \item Transmission power: 3 levels (low, medium, high)
    \item Distribution strategy: 2 options (single-path, multi-path)
    \item Packet size: 3 sizes (small, medium, large)
    \item Additional parameters depending on experiment setup
\end{itemize}

\textbf{Adaptation space scale}:
\begin{itemize}
    \item Small: 216 discrete options
    \item Medium: 1,296 discrete options
    \item Large: 4,096 discrete options
\end{itemize}

These scales represent realistic adaptation spaces in self-adaptive systems, challenging for both search-based and learning-based approaches.

\subsection{Implications for search/learning complexity}

\textbf{Search complexity}:
\begin{itemize}
    \item Exhaustive search requires evaluating all options
    \item Verification cost per option: 100s of simulations
    \item Total search time becomes prohibitive in large spaces
\end{itemize}

\textbf{Learning complexity}:
\begin{itemize}
    \item Large action spaces require more training experiences
    \item Risk of overfitting to training environments
    \item Need for sample-efficient learning algorithms
\end{itemize}

\section{RL Formalization for the Adaptive Problem (exact equations from paper)}

This section provides the formal MDP formulation for DeltaIoT adaptation.

\subsection{State space definition (e.g., $S = (E, P, L)$)}

\textbf{State representation}: $s_t = (E_t, P_t, L_t)$ where:
\begin{itemize}
    \item $E_t$: Energy consumption level
    \item $P_t$: Packet loss rate
    \item $L_t$: Latency measure
\end{itemize}

Additional state features may include:
\begin{itemize}
    \item Current configuration parameters
    \item Recent trend indicators
    \item Environmental observations (interference estimates)
\end{itemize}

State space dimensionality: $s \in \mathbb{R}^d$ where $d$ depends on feature engineering.

\subsection{Action space: discrete adaptation options and encoding}

\textbf{Action space}: $\mathcal{A} = \{a_1, a_2, \ldots, a_n\}$ where $n \in \{216, 1296, 4096\}$

\textbf{Action encoding}:
\begin{itemize}
    \item Each action $a_i$ corresponds to a unique configuration
    \item Integer encoding: $a_i \in \{0, 1, \ldots, n-1\}$
    \item One-hot encoding for neural network input (if needed)
\end{itemize}

\subsection{Reward functions: multi-objective normalization, threshold/setpoint formulations (include equations 11--12 and Table 2)}

\textbf{Multi-objective reward function}:

\begin{equation}
r_t = w_E \cdot r_E(s_t) + w_P \cdot r_P(s_t) + w_L \cdot r_L(s_t)
\end{equation}

where $w_E + w_P + w_L = 1$ are objective weights.

\textbf{Threshold-based formulation} (Equation 11 from article):

\begin{equation}
r_i(s_t) = \begin{cases}
1 & \text{if } m_i(s_t) \leq \theta_i \\
-1 & \text{otherwise}
\end{cases}
\end{equation}

where $m_i(s_t)$ is metric $i$ value and $\theta_i$ is threshold.

\textbf{Setpoint-based formulation} (Equation 12 from article):

\begin{equation}
r_i(s_t) = 1 - \frac{|m_i(s_t) - \theta_i|}{\theta_i}
\end{equation}

\textbf{Table 2}: Reward function parameters
\begin{table}[H]
\centering
\begin{tabular}{lll}
\textbf{Metric} & \textbf{Threshold $\theta_i$} & \textbf{Weight $w_i$} \\
\hline
Energy & 0.5 & 0.33 \\
Packet Loss & 0.1 & 0.34 \\
Latency & 100 ms & 0.33 \\
\end{tabular}
\caption{Reward function parameters for DeltaIoT experiments.}
\end{table}

\subsection{Transition dynamics (how verifier / ActivFORMS produces next states)}

\textbf{Transition function}: $s_{t+1} \sim P(s'|s_t, a_t)$

\textbf{Verification-based transitions}:
\begin{itemize}
    \item UPPAAL-SMC statistical model checker simulates system behavior
    \item ActivFORMS execution engine provides runtime verification
    \item Multiple simulation runs produce probability distributions
    \item Next state sampled from distribution: $s_{t+1} \sim P(\cdot|s_t, a_t)$
\end{itemize}

\textbf{Stochastic transitions}:
\begin{itemize}
    \item Same $(s_t, a_t)$ pair can produce different $s_{t+1}$
    \item Captures uncertainty in system behavior
    \item Requires learning robust policies
\end{itemize}

\section{Continuous / Streaming Training Assumptions}

This section clarifies the continuous learning paradigm used in RS-DRL.

\subsection{Treating adaptation cycles as a continuous stream (training horizons and validation splits per instance)}

\textbf{Continuous stream paradigm}:
\begin{itemize}
    \item Adaptation cycles arrive continuously during operation
    \item Each cycle generates experience tuple $(s_t, a_t, r_t, s_{t+1})$
    \item Experiences stored in replay memory for training
    \item No episodic boundaries in runtime operation
\end{itemize}

\textbf{Training horizons}:
\begin{itemize}
    \item Training occurs over fixed time horizons (e.g., 1000 cycles)
    \item Validation performed on separate environment instances
    \item Cross-validation across interference levels
\end{itemize}

\subsection{Temporal separation: runtime adaptation vs asynchronous learning (design assumption)}

\textbf{Key architectural assumption}: Runtime adaptation and learning are temporally separated.

\textbf{Runtime adaptation}:
\begin{itemize}
    \item Uses pre-trained model for inference
    \item Pure forward pass through neural network
    \item No gradient computation or weight updates
    \item Fast, deterministic decision-making
\end{itemize}

\textbf{Asynchronous learning}:
\begin{itemize}
    \item Occurs in background process
    \item Samples experiences from replay memory
    \item Updates model weights through gradient descent
    \item Stores updated model for future runtime use
\end{itemize}

This separation ensures runtime stability while enabling continuous improvement.

\section{Verification \& Runtime Models for Option Analysis}

Verification provides quality estimates for adaptation options under uncertainty.

\subsection{UPPAAL-SMC models and ActivFORMS execution engine (role in reward calculation)}

\textbf{UPPAAL-SMC}:
\begin{itemize}
    \item Statistical model checker for stochastic timed automata
    \item Simulates system behavior under uncertainty
    \item Produces probability estimates for property satisfaction
    \item Used during training to evaluate adaptation options
\end{itemize}

\textbf{ActivFORMS}:
\begin{itemize}
    \item Runtime verification engine
    \item Executes formal models with current state parameters
    \item Provides quality-of-service predictions
    \item Generates reward signals for learning
\end{itemize}

\textbf{Role in reward calculation}:
\begin{enumerate}
    \item Agent selects action $a_t$ in state $s_t$
    \item Verifier evaluates $(s_t, a_t)$ using formal models
    \item Simulation produces QoS metrics: $(E, P, L)$
    \item Reward computed from metrics: $r_t = f(E, P, L)$
    \item Next state extracted from simulation: $s_{t+1}$
\end{enumerate}

\section{Formal Problem Statement (thesis version) --- precise objective and constraints (ready to feed Chapter 4)}

\textbf{Formal Problem Statement}:

Given:
\begin{itemize}
    \item Self-adaptive system with discrete adaptation space $\mathcal{A}$, $|\mathcal{A}| \in [100, 10000]$
    \item State space $\mathcal{S} \subset \mathbb{R}^d$ representing system and environment
    \item Stochastic transition dynamics $P(s'|s,a)$ with runtime uncertainty
    \item Multi-objective reward function $R(s,a,s')$ encoding QoS goals
    \item Heterogeneous environment instances $\mathcal{E} = \{e_1, \ldots, e_k\}$
\end{itemize}

Find:
\begin{itemize}
    \item Policy $\pi^*: \mathcal{S} \rightarrow \mathcal{A}$ that maximizes expected cumulative reward:
    \begin{equation}
    \pi^* = \arg\max_\pi \mathbb{E}_{\tau \sim \pi, e \sim \mathcal{E}} \left[ \sum_{t=0}^T \gamma^t R(s_t, a_t, s_{t+1}) \right]
    \end{equation}
\end{itemize}

Subject to constraints:
\begin{itemize}
    \item Sample efficiency: Learn with limited training experiences
    \item Generalization: Policy performs well across all $e \in \mathcal{E}$
    \item Continuous operation: System remains operational during learning
    \item Computational feasibility: Training and inference within resource bounds
\end{itemize}

This formulation guides the RS-DRL architecture and algorithm design in Chapter 4.

\newpage

% ============================================================================
% CHAPTER 4: ARCHITECTURE AND RS-DRL METHOD
% ============================================================================

\setcounter{section}{0}
\renewcommand{\thesection}{4.\arabic{section}}
\renewcommand{\thesubsection}{4.\arabic{section}.\arabic{subsection}}
\renewcommand{\thesubsubsection}{4.\arabic{section}.\arabic{subsection}.\arabic{subsubsection}}

\section{Chapter overview \& architectural perspective (scope, process-oriented organization)}

This chapter presents the RS-DRL architecture and method, organized by process flow rather than component structure.

\textbf{Chapter scope}:
\begin{itemize}
    \item System-level architecture integrating RS-DRL with MAPE-K
    \item Nine-stage adaptation cycle from monitoring to execution
    \item Asynchronous learning process with reward reshaping
    \item Knowledge management and process coordination
\end{itemize}

\textbf{Process-oriented organization}:
\begin{itemize}
    \item Each section answers: ``What happens next, and why?''
    \item MAPE-K serves as explanatory lens, not primary structure
    \item Decomposed figures aligned with section boundaries
    \item Complete workflows provided in appendices
\end{itemize}

\section{System-level Self-Adaptation Loop (nine-stage continuous adaptation cycle)}

The RS-DRL-based self-adaptive system operates through a nine-stage continuous cycle.

\subsection{Stage 1: Runtime data collection (sensors $\rightarrow$ Monitor)}

\textbf{Process}: Monitor collects runtime data from managed system and environment.

\textbf{Data sources}:
\begin{itemize}
    \item System metrics: energy consumption, packet loss, latency
    \item Environment observations: interference estimates, workload
    \item Configuration state: current parameter values
\end{itemize}

\textbf{Output}: Raw sensor data ready for preprocessing.

\subsection{Stage 2: State forwarding to RS-DRL module (experience collection)}

\textbf{Process}: Validated state vector forwarded to RS-DRL module for experience collection.

\textbf{Experience tuple construction}:
\begin{itemize}
    \item Previous state $s_t$
    \item Action taken $a_t$
    \item Reward received $r_t$
    \item Current state $s_{t+1}$
\end{itemize}

\textbf{Storage}: Experience $(s_t, a_t, r_t, s_{t+1})$ stored in replay memory.

\subsection{Stage 3: Asynchronous model training and storage (replay $\rightarrow$ train $\rightarrow$ store)}

\textbf{Process}: Background learning process trains model using replay memory.

\textbf{Training steps}:
\begin{enumerate}
    \item Sample minibatch from replay memory
    \item Apply randomized reward reshaping
    \item Compute TD targets and loss
    \item Update Q-network weights via gradient descent
    \item Store updated model in Knowledge repository
\end{enumerate}

\textbf{Asynchronous execution}: Training occurs independently of runtime adaptation.

\subsection{Stage 4: State forwarding to Analyzer}

\textbf{Process}: State vector forwarded to Analyzer for goal evaluation.

\textbf{Analyzer input}:
\begin{itemize}
    \item Current state $s_t$
    \item Adaptation goals and thresholds
    \item Historical state trajectory
\end{itemize}

\subsection{Stage 5: Adaptation triggering (Analyzer detects violations)}

\textbf{Process}: Analyzer evaluates whether adaptation is needed.

\textbf{Triggering conditions}:
\begin{itemize}
    \item QoS metric exceeds threshold
    \item Trend indicates future violation
    \item Environment change detected
    \item Periodic adaptation schedule
\end{itemize}

\textbf{Output}: Adaptation trigger signal sent to Planner.

\subsection{Stage 6: Trained model retrieval}

\textbf{Process}: Planner retrieves latest trained model from Knowledge repository.

\textbf{Model selection}:
\begin{itemize}
    \item Latest successfully trained model
    \item Model validated on current environment type
    \item Fallback to previous model if training failed
\end{itemize}

\subsection{Stage 7: Adaptation option prediction (Adaptation Option Predictor)}

\textbf{Process}: Planner uses model to predict best adaptation option.

\textbf{Prediction steps}:
\begin{enumerate}
    \item Forward pass: $Q(s_t, \cdot) = \text{model}(s_t)$
    \item Compute Q-values for all actions
    \item Select action: $a_t = \arg\max_a Q(s_t, a)$ (greedy) or $\epsilon$-greedy
\end{enumerate}

\textbf{Output}: Predicted adaptation option $a_t$.

\subsection{Stage 8: Adaptation option retrieval \& verification}

\textbf{Process}: Executor retrieves detailed specification and verifies option.

\textbf{Verification}:
\begin{itemize}
    \item UPPAAL-SMC simulates $(s_t, a_t)$ configuration
    \item Estimates QoS metrics and goal satisfaction probability
    \item Provides decision justification
\end{itemize}

\textbf{Output}: Verified adaptation specification ready for execution.

\subsection{Stage 9: Adaptation application (Executor)}

\textbf{Process}: Executor applies adaptation to managed system.

\textbf{Execution steps}:
\begin{enumerate}
    \item Translate adaptation option to system commands
    \item Apply configuration changes to managed system
    \item Verify successful reconfiguration
    \item Log adaptation event
\end{enumerate}

\textbf{Loop closure}: Monitoring resumes (Stage 1), closing the adaptation loop.

\section{Mapping RS-DRL to MAPE-K (conceptual mapping and roles)}

This section provides conceptual mapping of RS-DRL stages to MAPE-K components.

\textbf{MAPE-K mapping}:
\begin{itemize}
    \item \textbf{Monitor} (Stages 1, 2, 4): Data collection, state construction, dissemination
    \item \textbf{Analyze} (Stages 4, 5): Goal evaluation, adaptation triggering
    \item \textbf{Plan} (Stages 5, 6, 7): Model retrieval, option prediction
    \item \textbf{Execute} (Stages 8, 9): Verification, adaptation application
    \item \textbf{Knowledge} (Stages 3, 6, 8): Model storage, option repository, state history
\end{itemize}

\textbf{RS-DRL role}: Enhances Planning phase with learned intelligence, operates asynchronously in Knowledge component.

\section{Runtime Monitoring and State Construction}

Monitoring transforms raw sensor data into validated state representations.

\subsection{Data collection, preprocessing, feature extraction, state vector construction}

\textbf{Data collection}:
\begin{itemize}
    \item Sensors provide raw measurements
    \item Sampling rate: continuous or periodic
    \item Data buffering for temporal features
\end{itemize}

\textbf{Preprocessing}:
\begin{itemize}
    \item Filtering: Remove noise and outliers
    \item Normalization: Scale features to $[0,1]$ or standardize
    \item Aggregation: Compute moving averages, trends
\end{itemize}

\textbf{Feature extraction}:
\begin{itemize}
    \item QoS metrics: energy, packet loss, latency
    \item Derived features: rates of change, violation indicators
    \item Context features: current configuration, time of day
\end{itemize}

\textbf{State vector construction}:
\begin{itemize}
    \item Assemble features into vector $s_t \in \mathbb{R}^d$
    \item Validate: Check for missing values, range violations
    \item Disseminate: Forward to Analyzer and RS-DRL module
\end{itemize}

\section{Analysis and Adaptation Triggering}

Analysis evaluates system state against adaptation goals and triggers planning when needed.

\subsection{Change detection \& thresholding assumptions (when retraining is triggered)}

\textbf{Goal evaluation}:
\begin{itemize}
    \item Compare metrics against thresholds: $m_i(s_t) \leq \theta_i$
    \item Evaluate multi-objective satisfaction
    \item Detect trends indicating future violations
\end{itemize}

\textbf{Adaptation triggering}:
\begin{itemize}
    \item Immediate: Threshold violation detected
    \item Predictive: Trend analysis predicts violation
    \item Periodic: Time-based adaptation schedule
\end{itemize}

\textbf{Retraining triggers}:
\begin{itemize}
    \item Environment change detected (interference shift)
    \item Policy performance degradation
    \item Sufficient new experiences collected
    \item Scheduled retraining interval
\end{itemize}

\section{Planning \& Adaptation Option Prediction}

Planning uses the trained RS-DRL model to select adaptation options.

\subsection{Adaptation Option Selector ($\epsilon$-greedy + DQN) \& architecture of predictor}

\textbf{Adaptation Option Predictor architecture}:
\begin{itemize}
    \item Input layer: State vector $s_t \in \mathbb{R}^d$
    \item Hidden layers: Fully connected neural network
    \item Output layer: Q-values for all actions $Q(s_t, a_i)$, $i = 1, \ldots, n$
\end{itemize}

\textbf{Action selection}:
\begin{itemize}
    \item Training: $\epsilon$-greedy exploration
    \begin{equation}
    a_t = \begin{cases}
    \text{random action} & \text{with probability } \epsilon \\
    \arg\max_a Q(s_t, a) & \text{with probability } 1-\epsilon
    \end{cases}
    \end{equation}
    \item Runtime: Greedy selection ($\epsilon = 0$)
    \begin{equation}
    a_t = \arg\max_a Q(s_t, a)
    \end{equation}
\end{itemize}

\textbf{Pure inference at runtime}: No learning, no replay sampling, no weight updates.

\section{Verification of Adaptation Options}

Verification provides quality estimates and decision justification for selected options.

\subsection{Runtime model checking with UPPAAL-SMC \& how verifier informs reward}

\textbf{Verification process}:
\begin{enumerate}
    \item Parameterize runtime model with $(s_t, a_t)$
    \item Execute UPPAAL-SMC statistical model checking
    \item Run multiple simulations (e.g., 1000 runs)
    \item Collect QoS metrics from simulations
    \item Compute quality estimates and probabilities
\end{enumerate}

\textbf{Dual role of verification}:
\begin{itemize}
    \item \textbf{Runtime}: Provides decision justification and confidence estimates
    \item \textbf{Learning}: Generates reward signals $r_t$ from QoS metrics
\end{itemize}

\textbf{Reward generation}:
\begin{equation}
r_t = f(\text{QoS metrics from verification})
\end{equation}

where $f$ is the multi-objective reward function defined in Section 3.4.3.

\section{Execution and Adaptation Application}

Execution translates selected adaptation options into system reconfigurations.

\subsection{Executor behavior and coordination with Knowledge repository}

\textbf{Adaptation specification retrieval}:
\begin{itemize}
    \item Query Knowledge repository with action ID $a_t$
    \item Retrieve complete configuration specification
    \item Includes all parameter values and dependencies
\end{itemize}

\textbf{Execution process}:
\begin{enumerate}
    \item Validate adaptation specification
    \item Translate to system-specific commands
    \item Apply configuration changes atomically
    \item Verify successful reconfiguration
    \item Update system state in Knowledge
\end{enumerate}

\textbf{Error handling}:
\begin{itemize}
    \item Rollback on execution failure
    \item Fallback to safe default configuration
    \item Log error for analysis
\end{itemize}

\section{Asynchronous RS-DRL Learning and Policy Update (core algorithmic chapter)}

This section presents the core RS-DRL algorithm with randomized reward reshaping.

\subsection{Overall RS-DRL algorithm (Algorithm 1) and parameters ($\eta$, $\gamma$, $\epsilon$, $\rho$)}

\textbf{Algorithm 1: RS-DRL Training}

\textbf{Parameters}:
\begin{itemize}
    \item $\eta$: Learning rate
    \item $\gamma$: Discount factor
    \item $\epsilon$: Exploration rate
    \item $\rho$: Reward reshaping probability
\end{itemize}

\textbf{Algorithm outline}:
\begin{enumerate}
    \item Initialize Q-network $Q(s,a;\theta)$ and target network $Q(s,a;\theta^-)$
    \item Initialize replay memory $\mathcal{D}$
    \item For each training episode:
    \begin{enumerate}
        \item Observe initial state $s_0$
        \item For each time step $t$:
        \begin{enumerate}
            \item Select action $a_t$ using $\epsilon$-greedy policy
            \item Execute action, observe reward $r_t$ and next state $s_{t+1}$
            \item Store experience $(s_t, a_t, r_t, s_{t+1})$ in $\mathcal{D}$
            \item Sample minibatch from $\mathcal{D}$
            \item Apply randomized reward reshaping (Algorithm 2)
            \item Compute TD targets and update Q-network
        \end{enumerate}
        \item Periodically update target network: $\theta^- \leftarrow \theta$
    \end{enumerate}
    \item Store trained model in Knowledge repository
\end{enumerate}

\subsection{Replay Memory \& randomized reward reshaping (Algorithm 2) --- rationale \& implementation}

\textbf{Algorithm 2: Randomized Reward Reshaping}

\textbf{Input}: Minibatch of experiences $\{(s_i, a_i, r_i, s'_i)\}$

\textbf{Output}: Reshaped minibatch $\{(s_i, a_i, \tilde{r}_i, s'_i)\}$

\textbf{Algorithm}:
\begin{enumerate}
    \item For each experience $(s_i, a_i, r_i, s'_i)$ in minibatch:
    \begin{enumerate}
        \item If $r_i < 0$ (failed experience):
        \begin{enumerate}
            \item With probability $\rho$: $\tilde{r}_i = |r_i|$ (reshape to positive)
            \item With probability $1-\rho$: $\tilde{r}_i = r_i$ (keep original)
        \end{enumerate}
        \item Else: $\tilde{r}_i = r_i$ (keep positive rewards unchanged)
    \end{enumerate}
    \item Return reshaped minibatch
\end{enumerate}

\textbf{Rationale}:
\begin{itemize}
    \item Treats certain failures as optimistic successes
    \item Provides additional positive learning signals
    \item Accelerates exploration of promising regions
    \item Promotes generalization across environments
    \item Randomization prevents overfitting to reshaped distribution
\end{itemize}

\textbf{Implementation details}:
\begin{itemize}
    \item Reshaping applied during training only, not at runtime
    \item Probability $\rho$ is hyperparameter (typically 0.1-0.3)
    \item Original experiences preserved in replay memory
    \item Reshaping creates temporary training signal
\end{itemize}

\subsection{Model training, target networks, loss function, stability practices}

\textbf{TD target computation}:
\begin{equation}
y_i = r_i + \gamma \max_{a'} Q(s'_i, a'; \theta^-)
\end{equation}

\textbf{Loss function}:
\begin{equation}
\mathcal{L}(\theta) = \mathbb{E}_{(s,a,r,s') \sim \mathcal{D}} \left[ (y - Q(s,a;\theta))^2 \right]
\end{equation}

\textbf{Gradient update}:
\begin{equation}
\theta \leftarrow \theta - \eta \nabla_\theta \mathcal{L}(\theta)
\end{equation}

\textbf{Target network update}:
\begin{itemize}
    \item Periodic hard update: $\theta^- \leftarrow \theta$ every $C$ steps
    \item Or soft update: $\theta^- \leftarrow \tau \theta + (1-\tau) \theta^-$
\end{itemize}

\textbf{Stability practices}:
\begin{itemize}
    \item Gradient clipping: Prevent exploding gradients
    \item Reward normalization: Stabilize learning across scales
    \item Experience replay: Break temporal correlations
    \item Target network: Stabilize TD targets
    \item Learning rate scheduling: Reduce $\eta$ over time
\end{itemize}

\section{Knowledge Management and Process Coordination}

Knowledge serves as shared memory coordinating heterogeneous processes.

\subsection{Knowledge repository design (models, adaptation options, state history)}

\textbf{Knowledge components}:
\begin{itemize}
    \item \textbf{Model repository}: Trained Q-networks, metadata, versioning
    \item \textbf{Adaptation option catalog}: Configuration specifications, parameters
    \item \textbf{State history}: Time-series of system states and adaptations
    \item \textbf{Replay memory}: Experience tuples for training
\end{itemize}

\textbf{Data structures}:
\begin{itemize}
    \item Models: Serialized neural network weights, architecture specs
    \item Options: JSON/XML configuration files
    \item History: Time-series database or log files
    \item Replay: Circular buffer or database
\end{itemize}

\subsection{Read/write coordination contracts and heterogeneous process synchronization}

\textbf{Read/write contracts}:
\begin{itemize}
    \item \textbf{Monitor writes}: State vectors (Stage 1)
    \item \textbf{Learning writes}: Trained models (Stage 3)
    \item \textbf{Planner reads}: Trained models (Stage 6)
    \item \textbf{Executor reads}: Adaptation specifications (Stage 8)
\end{itemize}

\textbf{Synchronization mechanisms}:
\begin{itemize}
    \item Atomic writes: Ensure model updates are complete
    \item Versioning: Track model versions, enable rollback
    \item Locks: Prevent concurrent write conflicts
    \item Notifications: Signal when new models available
\end{itemize}

\textbf{Temporal decoupling}:
\begin{itemize}
    \item Runtime processes read latest available model
    \item Learning processes write asynchronously
    \item No blocking between runtime and learning
\end{itemize}

\section{Implementation details \& hyperparameter choices}

This section provides practical implementation details for reproducibility.

\subsection{Gymnasium environment, Stable-Baselines3 DQN extension, dataset traces (Gheibi \& Weyns)}

\textbf{Gymnasium environment}:
\begin{itemize}
    \item Custom environment wrapping DeltaIoT simulator
    \item State space: $\mathbb{R}^d$ (continuous)
    \item Action space: Discrete($n$) where $n \in \{216, 1296, 4096\}$
    \item Reward: Multi-objective function from Section 3.4.3
    \item Episode termination: Fixed horizon or goal achievement
\end{itemize}

\textbf{Stable-Baselines3 DQN extension}:
\begin{itemize}
    \item Base: SB3 DQN implementation
    \item Extension: Custom callback for reward reshaping
    \item Integration: Modify replay buffer sampling
    \item Logging: Track reshaping statistics
\end{itemize}

\textbf{Dataset traces}:
\begin{itemize}
    \item Source: Gheibi \& Weyns DeltaIoT datasets
    \item Interference levels: 0.0, 0.2, 0.4, 0.6, 0.8, 1.0
    \item Message loads: Low, medium, high
    \item Training/validation split: 70/30 or cross-validation
\end{itemize}

\textbf{Hyperparameter choices}:
\begin{table}[H]
\centering
\begin{tabular}{ll}
\textbf{Parameter} & \textbf{Value} \\
\hline
Learning rate $\eta$ & 0.0001 \\
Discount factor $\gamma$ & 0.99 \\
Exploration rate $\epsilon$ & 0.1 (decayed) \\
Reshaping probability $\rho$ & 0.2 \\
Replay buffer size & 100,000 \\
Minibatch size & 64 \\
Target update frequency & 1,000 steps \\
Network architecture & [128, 128] \\
\end{tabular}
\caption{RS-DRL hyperparameter settings.}
\end{table}

\section{Chapter summary and transition to evaluation (Chapter 5)}

This chapter presented the RS-DRL architecture and method:

\textbf{Key contributions}:
\begin{itemize}
    \item Nine-stage adaptation cycle integrating RS-DRL with MAPE-K
    \item Temporal separation of runtime adaptation and asynchronous learning
    \item Randomized reward reshaping algorithm for accelerated learning
    \item Knowledge-based coordination of heterogeneous processes
    \item Practical implementation using Gymnasium and Stable-Baselines3
\end{itemize}

\textbf{Architectural principles}:
\begin{itemize}
    \item Process-oriented organization for traceability
    \item Verification as justification for learning-based decisions
    \item Asynchronous learning for continuous operation
    \item Modular design enabling independent component evolution
\end{itemize}

\textbf{Transition to evaluation}:

Chapter 5 evaluates RS-DRL through comprehensive experiments on DeltaIoT:
\begin{itemize}
    \item Learning performance: Convergence speed, sample efficiency
    \item Generalization: Cross-environment policy transfer
    \item Comparison: Baseline DQN, HER, search-based methods
    \item Ablation studies: Impact of reshaping probability $\rho$
    \item Scalability: Performance across adaptation space sizes
\end{itemize}

\end{document}
